% The collective evidence from recent literature indicates that DRL is progressively transforming PECs research from model-dependent analysis toward data-driven intelligence. Across design, control, and maintenance stages, DRL has demonstrated the capacity to uncover cross-domain relationships that are often inaccessible to traditional analytical or optimization frameworks. This paradigm shift suggests that converter systems can evolve from manually tuned and locally optimized devices to self-adaptive and globally coordinated energy conversion platforms.

% A key insight emerging from current studies is that the effectiveness of DRL in PECs does not primarily stem from algorithmic novelty, but from its integration with physical knowledge, hierarchical structures, and hardware constraints. Physics-informed reward design, surrogate-assisted learning, and hybrid DRL–model-based architectures consistently yield faster convergence, greater safety, and improved interpretability. These findings indicate that embedding domain principles into DRL pipelines is essential for bridging the gap between theoretical learning performance and practical converter operation.

% In converter control, DRL has shown particular promise in coordinating multi-objective trade-offs—such as efficiency, stability, and dynamic response—without explicit modeling of complex nonlinearities. The literature reveals that hybrid DRL control, which couples autonomous exploration with deterministic control priors, achieves a favorable balance between adaptability and robustness. This convergence of learning and control theory points toward a future in which DRL serves as an adaptive supervisory layer within safety-certified hierarchical control architectures.

% Within the maintenance domain, DRL-driven approaches to condition monitoring, fault diagnosis, and RUL prediction provide a blueprint for realizing predictive and self-healing converter systems. The combination of reinforcement learning with attention-based temporal modeling and transfer learning offers scalable solutions for real-time health management under heterogeneous conditions. These developments mark an evolution from periodic inspection toward continuous, data-informed reliability engineering.

% Looking ahead, several research priorities emerge. Future work should focus on unifying DRL with physical modeling, uncertainty quantification, and embedded implementation. Hardware-aware policy deployment, multi-agent coordination for modular converters, and continual learning for long-term adaptability are expected to define the next phase of advancement. Equally important is the establishment of standardized validation protocols, interpretable decision frameworks, and safety certification pathways to ensure trust and reproducibility.

% In summary, DRL provides a coherent foundation for developing autonomous, resilient, and optimally coordinated power electronic systems. Its future impact will depend less on algorithmic complexity and more on the ability to integrate learning intelligence with physical realism and engineering practicality. Achieving this synthesis represents the key direction for translating DRL from an experimental tool into an industrially reliable technology for next-generation power electronics.

DRL is redefining research in PECs not merely through algorithmic novelty, but more importantly through the systemic integration of data-driven intelligence with physical modeling and hardware-aware design methodologies. Across the design, control, and maintenance stages, DRL has evolved from an exploratory optimization tool into a cohesive framework capable of learning, adapting, and coordinating converter behavior under nonlinear and time-varying conditions. This transformation marks a broader paradigm shift from model-dependent analysis to knowledge-guided autonomy in power electronics.

Within converter \textit{design}, DRL enables topology exploration, parameter optimization, and multi-objective trade-off analysis through adaptive, data-driven reasoning that complements analytical constraints. In \textit{control}, hybrid DRL architectures have demonstrated the ability to reconcile adaptability with robustness by integrating autonomous policy learning with deterministic control priors. In \textit{maintenance}, DRL-based methods for condition monitoring, anomaly detection, and RUL prediction have accelerated the transition from periodic inspection to continuous, predictive reliability management. Collectively, these developments illustrate the expanding role of DRL as a unifying intelligence layer across the converter life cycle.

A central insight emerging from the literature is that the practical effectiveness of DRL in PECs depends less on algorithmic complexity and more on its ability to incorporate domain knowledge, hierarchical structures, and physical constraints. Physics-informed reward shaping, surrogate-assisted policy learning, and model-based hybridization consistently enhance stability, convergence, and interpretability. These strategies bridge the persistent gap between theoretical reinforcement learning performance and the reliability standards required for industrial converter operation.

Looking forward, the next stage of advancement will depend on unifying learning intelligence with engineering realism. Key research priorities include: (i) safety-guaranteed and physics-constrained DRL for verifiable converter operation, (ii) hardware–algorithm co-design enabling real-time inference on FPGA, DSP, and SoC platforms, (iii) federated and multi-agent coordination for modular and distributed systems, and (iv) continual learning frameworks that preserve adaptability under aging, disturbance, and environmental variability. Equally important is the development of standardized benchmarks, explainable decision mechanisms, and certification pathways to ensure reproducibility and industrial trust.

In summary, DRL provides a coherent foundation for the development of autonomous, resilient, and optimally coordinated power electronic systems. Its ultimate maturity will be defined not by algorithmic sophistication but by the seamless integration of learning-based intelligence with physical laws, safety assurance, and embedded feasibility. Achieving this synthesis will mark the decisive step in transforming DRL from an experimental research paradigm into a dependable technology for next-generation PECs.
